\subsection{Primitive path slopes with arbitrary unit arrows}

\begin{theorem}[DP46: unit-arrow and invariant-frame independence]\label{DP46}
Let $p$ be odd, $h\ge3$, $R=W(k)$ with perfect $k\supset\F_{p^h}$.
Take DP40's singleton cyclic display with arbitrary units $u_i$:
\[
 Fe_0=u_0e_1,\qquad Fe_i=pu_ie_{i+1}\quad(1\le i\le d=h-1),
\]
and its actual complementary dual. Put $a_i=w_0(u_i)\in R^\times$.
Use its Teichm\"uller coordinates and the invariant primitive Poincar\'e frame.
Every conclusion of DP42 about the exact primitive mixed-line threshold, finite norm, factorial bound and liminf remains valid, with
$k_*=h-j$ the reduced dual direction's restricted-$p$ nilpotency length.
The exact coefficient congruence is now
\[
 p^{S_{k_*}}[t^{p^{k_*}+1}]\beta(ut,\mathbf vt)
 \equiv(-1)^{k_*-1}\theta_j u v_j^{p^{k_*}}\pmod p,\qquad
 \theta_j=\prod_{r=j}^{d}a_r^{-p^{d-r}}\in R^\times.
\]
The same assertions hold for every normalized invariant integral frame with its invariant logarithmic splitting, for every unit exponential parameter $\lambda\in R^\times$.
They also transport under supplied integral CM isomorphisms over unramified extensions. No finite unramified trivialization of arbitrary twists is assumed.
\end{theorem}

\begin{proof}
Let $B_n$ be the $n$th ghost matrix of $F$, and $B'_n$ its complementary reciprocal-unit counterpart. Thus
$B_n^{\mathsf t}B'_n=pI$, and every ghost of each $u_i$ is a unit.
Set $P_{n,0}=I$ and
$P_{n,r}=B_nB_{n+1}\cdots B_{n+r-1}$, with analogous $P'$.
On divided ghosts the resolution is $z_n\mapsto B_nz_{n+1}-z_n$, whose inverse on the finite rational sections is
$-\sum_{r\ge0}P_{n,r}z_{n+r}$.
The logarithm is the selected projection of $\sum_nP_{0,n}z_n$.
DP40's uniform ghost-denominator bound and positive $p$-power per complete cycle justify these sums on the completed integral quotient.

For the Teichm\"uller sections put
$\xi_n=S^{p^n}e_0/p^n$ and
$\upsilon_n=\sum_iY_i^{p^n}e_i^*/p^n$.
Let $C_H$ be the $e_0^*$ component of $\sum_nP'_{0,n}\upsilon_n$.
The additive rational sections with only zeroth ghosts
$L_G(S)e_0$ and $\sum_i(L_H)_i e_i^*$ have relation-module defects
\[
 -\sum_rP_{n,r}\xi_{n+r}+\mathbf1_{n=0}L_G(S)e_0,\qquad
 -\sum_rP'_{n,r}\upsilon_{n+r}
       +\mathbf1_{n=0}\sum_i(L_H)_ie_i^*.
\]
Their selected zeroth components vanish. Substitution in the exact quotient-frame comparison of DP41, using the same negative $\operatorname{hex}$ convention, gives
\begin{equation}\label{eq:unitprimitive}
 b(S,\mathbf Y)=
 -\sum_{n,r\ge0}p^n
       \langle P_{n,r+1}\xi_{n+r+1},\upsilon_n\rangle
       +L_G(S)C_H(\mathbf Y).
\end{equation}
This uses no commutation between distinct ghost matrices.

Unique cyclic paths show that the first sum, the line logarithm, and the complementary component have the exact forms
\[
 -\sum_{i=1}^d\sum_{l,n\ge0}
   \eta_{i,l,n}\frac{S^{p^{n+lh+i}}Y_i^{p^n}}{p^{n+l+1}},\qquad
 L_G=\sum_{a\ge0}A_aS^{q^a}/p^a,
\]
\[
 C_H=\sum_{i=1}^d\sum_{l\ge0}
        C_{i,l}Y_i^{p^{lh+h-i}}/p^{dl+h-i},
\]
where every $\eta,A_a,C_{i,l}$ is a unit and $A_0=1$.
A line path of length $lh+i$ contributes $dl+i-1$ numerator powers of $p$; a dual path to index zero first has length $h-i$ and gains one numerator $p$ per additional cycle. These counts prove the denominators, not just lower bounds for them.
The unique coefficient of $SY_d^p$ is $a_d^{-1}/p$.
The ordinary first dual-logarithm arrows are
$a_{i-1}^{-1}Y_{i-1}^p/p$ into coordinate $i\ge2$.

Use exactly DP42's interval and weights
$\omega_0=w$, $\omega_i=s-p^{i-j}(s-w)$.
Choose all scales $\alpha_i=\pi^{e\omega_i}$ using a common $\pi^e=p$, with $e$ clearing the weights, and put $\gamma=\alpha_0=\alpha_j$.
The supports and denominator valuations just proved are precisely those used in its strict initial-form inequalities. Its first-family baseline gap
$(s-w)(p^{k_*}-p+1-p^{1-j})$ is positive; its product-family gap
$(p^{h-i}-1)/(p-1)-(h-i)$ vanishes only for $i=d$; additional cycles strictly increase valuation. Thus
\[
 \overline{\mu^{-1}b(\alpha_0S,\alpha_1Y_1,\ldots,\alpha_dY_d)}
       =\bar a_d^{-1}SY_d^p,\qquad \mu=\alpha_0\alpha_d^p/p.
\]
The reduced contracted logarithm has ordinary arrows
$Y_i+\bar a_{i-1}^{-1}Y_{i-1}^p$, with the exceptional arrow absent.
Only the first unit input $j$ survives, exactly as in DP42.
Recursive inversion gives its last coordinate
\[
 (-1)^{d-j}
 \left(\prod_{r=j}^{d-1}\bar a_r^{-p^{d-r-1}}\right)
 \bar v_j^{p^{d-j}}T^{p^{d-j}}.
\]
Hence the entire restricted initial form is
$(-1)^{k_*-1}\bar\theta_j\bar u\bar v_j^{p^{k_*}}T^{p^{k_*}+1}$.
It has nonzero derivative. The identity
$\mu=\gamma^{p^{k_*}+1}/p^{S_{k_*}}$ gives the coefficient congruence.
DP41's nonzero-multiplier lemma below the threshold, the factorial estimate at the threshold, monotonicity, and Gauss multiplicativity now prove every asserted norm and factorial conclusion.
The first dual arrows also show that the restricted operation takes each $D_i$ to a unit times $D_{i+1}$ and kills the last. Its length on the primitive vector is $h-j$, as asserted.

The display torus $e_i\mapsto[\zeta^{p^i}]e_i$ commutes with all unit arrows and acts on the paired coordinates as
$(S,Y_i)\mapsto(\zeta S,\zeta^{-p^i}Y_i)$.
The pairing frame and additive sections are equivariant. No bilinear additive monomial $xy_i$ is invariant, so the invariant logarithmic splitting is unique in a fixed invariant frame.
If another such integral frame differs by an invariant normalized unit $J$, its splitting changes by $\pm\log J$.
For an invariant nonconstant monomial $S^a\mathbf Y^{\mathbf b}$ its original contraction weight is
\[
 \rho\left(a-\sum_ip^ib_i\right)+s|\mathbf b|.
\]
The parenthesis is divisible by $q-1$; this positive weight lies in $s\Z$ and is at least $s$.
Each DP42 profile weight strictly exceeds its original contraction weight. Uniform strictness follows from the minimum of these finitely many positive increments. Therefore the contracted $J-1$, and then $\log J$, have coefficient infimum strictly greater than $s$.
For every integral $\lambda$, $\exp(\pm\lambda\log J)$ is an integral unit there. It leaves both the threshold and the initial coefficient congruence unchanged; monotonicity handles larger and smaller scalar path weights.
This argument does not assert unrestricted binomial integrality for arbitrary $\lambda\in R$. For arbitrary, possibly noninvariant integral frames, integral binomial powers give the corresponding frame conclusion when $\lambda\in\Z_p$.

Finally a supplied integral $O_q$-linear isomorphism $\phi:G\to G_0$ and its inverse dual $\psi:H\to H_0$ transport the primitive pairing. Their tangent maps are a unit scalar $a$ and an invertible integral matrix $B$. Normalized logarithms satisfy
$L_{G_0}\phi=aL_G$ and $L_{H_0}\psi=BL_H$.
Primitive projections remain primitive, restricted-$p$ length is unchanged by functoriality, and integral prefactors and nonzero restrictions are preserved. Invariant frames compare by the unit just treated. Unramified extension does not change coefficient valuations, so numerical conclusions descend whenever the original expression is defined over the smaller ring.
\end{proof}

